%----------------------------------------------------------
% 제5장. 산업별 AI 거버넌스 적용
%----------------------------------------------------------

\chapter{산업별 AI 거버넌스 적용}
\label{ch:industry_governance}

AI 거버넌스의 원칙과 프레임워크는 보편적이지만, 
그 적용 방식은 산업의 특성, 규제 환경, 리스크 프로파일에 따라 달라져야 한다. 
본 장에서는 AI 활용이 활발하고 거버넌스 요구가 높은 
네 개 산업---금융, 의료, 공공, 제조---에서의 
AI 거버넌스 적용 방법과 고려사항을 다룬다.

각 산업별로 규제 환경, 주요 AI 활용 사례, 
특수한 윤리적 고려사항, 거버넌스 적용 전략을 분석하고, 
실무에서 활용할 수 있는 산업별 체크리스트를 제공한다.


%----------------------------------------------------------
\section{금융 산업}
\label{sec:5.1}

\subsection{금융 AI의 규제 환경}
\label{sec:5.1.1}

금융 산업은 AI 거버넌스 측면에서 가장 엄격한 규제를 받는 분야 중 하나이다. 
금융 AI는 개인의 경제적 기회에 직접적인 영향을 미치므로, 
Fairness과 투명성에 대한 요구가 특히 높다.

\paragraph{주요 규제 및 Guide라인}

\begin{table}[htbp]
\centering
\caption{금융 AI 관련 주요 규제}
\label{tab:financial_regulations}
\begin{tabularx}{\textwidth}{p{3.5cm}p{3cm}X}
\toprule
\textbf{규제/Guide라인} & \textbf{관할} & \textbf{AI 관련 핵심 내용} \\
\midrule
신용정보의 이용 및 보호에 관한 법률 & 
한국 & 
신용평가 시 성별, 출신 지역 등 사용 제한; Automation된 평가에 대한 설명 요구권 (제36조의2) \\
\midrule
금융분야 AI Guide라인 (금융위원회, 2021) & 
한국 & 
AI 활용 원칙, 내부제어, 소비자 보호, 설명가능성 요구 \\
\midrule
AI 기본법 시행령 & 
한국 & 
대출.신용평가를 고위험 AI로 분류 (제10조); 공정성.투명성 요구 \\
\midrule
EU AI Act & 
EU & 
신용평가, 보험 리스크 평가를 고위험 AI로 분류 (Annex III) \\
\midrule
SR 11-7 (Model 리스크 관리) & 
미국 (Fed/OCC) & 
모델 리스크 관리, 독립적 검증, 문서화 요구 \\
\midrule
Fair Lending Laws (ECOA, FHA) & 
미국 & 
대출 결정에서 차별 금지, Adverse Action Notice 요구 \\
\bottomrule
\end{tabularx}
\end{table}


\paragraph{금융위원회 AI Guide라인 핵심 Principle}

금융위원회의 ``금융분야 AI(AI) Guide라인''(2021)은 
다음 다섯 가지 원칙을 제시한다:

\begin{enumerate}
 \item \textbf{공정성}: AI 알고리즘이 특정 집단에 불합리한 차별을 초래하지 않도록 관리
 \item \textbf{설명가능성}: 소비자가 AI based 결정 주요 요인을 이해할 수 있도록 설명 제공
 \item \textbf{책임성}: AI 활용에 대한 명확한 책임 체계 구축
 \item \textbf{안전성}: AI 시스템의 보안 및 안정적 운영 보장
 \item \textbf{투명성}: AI 활용 사실 및 범위에 대한 고객 고지
\end{enumerate}


\subsection{금융 AI 주요 활용 사례와 리스크}
\label{sec:5.1.2}

\begin{table}[htbp]
\centering
\caption{금융 AI 활용 사례별 리스크 분석}
\label{tab:financial_ai_risks}
\begin{tabularx}{\textwidth}{p{2.5cm}p{3cm}p{3cm}X}
\toprule
\textbf{활용 Case} & \textbf{주요 리스크} & \textbf{영항받는 집단} & \textbf{거버넌스 중점 사항} \\
\midrule
신용평가 & 
역사적 차별 재생산, 프록시 차별 & 
저소득층, 특정 지역 거주자, 청년층 & 
공정성 평가 필수, 대안 데이터 사용 시 주의, 설명 제공 \\
\midrule
대출 심사 & 
차별적 승인/거절, 금리 Differential & 
소수 집단, 여성, 고령자 & 
4/5 규칙 적용, Adverse Action 설명, 인간 검토 \\
\midrule
보험 인수/가격 & 
건강 상태, 유전 정보 based 차별 & 
만성질환자, 장애인 & 
금지 변수 명확화, 가격 Differential 근거 투명화 \\
\midrule
이상거래 탐지 (FDS) & 
특정 집단 과탐지, 정상 거래 차단 & 
해외 거주자, 특정 업종 종사자 & 
FPR 집단별 분석, 이의 제기 절차 \\
\midrule
투자 자문/로보어드바이저 & 
부적합 추천, 집중 리스크 & 
금융 지식 부족 고객 & 
적합성 평가, 리스크 고지, 인간 감독 \\
\midrule
고객 서비스 (챗봇) & 
오정보, 취약 고객 Response 미흡 & 
고령자, 장애인 & 
에스컬레이션 절차, 정확성 모니터링 \\
\bottomrule
\end{tabularx}
\end{table}


\subsection{금융 AI 공정성 평가}
\label{sec:5.1.3}

금융 AI에서 공정성 평가는 특히 중요하며, 
규제 기관의 검사 대상이 될 수 있다.

\paragraph{신용평가 모델 공정성 평가 프레임워크}

\begin{lstlisting}[language=Python, caption={Financial AI Fairness Evaluation}, label={lst:financial_fairness}]
class FinancialAIFairnessEvaluator:
    """Evaluates 4/5 rule and Adverse Impact for Financial AI."""
    def __init__(self, protected_attrs=['gender', 'age', 'region']):
        self.protected_attrs = protected_attrs

    def evaluate_4_5_rule(self, y_pred, sensitive_features):
        results = {}
        for attr in self.protected_attrs:
            groups = sensitive_features[attr].unique()
            rates = {g: y_pred[sensitive_features[attr]==g].mean() for g in groups}
            max_r = max(rates.values())
            # Adverse Impact Ratio (AIR)
            results[attr] = {g: r/max_r for g, r in rates.items()}
        return results

    def generate_adverse_action_notice(self, shap_vals, feat_desc):
        # Sort factors by negative impact
        sorted_neg = sorted(shap_vals.items(), key=lambda x: x[1])[:4]
        notice = "Credit Decision Notice\n\nKey Factors:\n"
        for i, (f, v) in enumerate(sorted_neg, 1):
            notice += f"{i}. {feat_desc.get(f, f)}\n"
        return notice
\end{lstlisting}


\subsection{금융 AI 거버넌스 체크리스트}
\label{sec:5.1.4}

\begin{table}[htbp]
\centering
\caption{[실무 도구 5-1] 금융 AI 거버넌스 체크리스트}
\label{tab:financial_checklist}
\begin{tabularx}{\textwidth}{|p{0.5cm}|X|p{1.5cm}|p{3cm}|}
\hline
\textbf{No.} & \textbf{Check 항목} & \textbf{결과} & \textbf{근거 규정} \\
\hline
\multicolumn{4}{|c|}{\textbf{1. 사전 승인}} \\
\hline
1.1 & AI 모델 사용 목적 및 범위 정의 & \checkmark / \times & 금융위 Guide라인 \\
\hline
1.2 & 리스크 평가 수행 (MRM 관점) & \checkmark / \times & SR 11-7 \\
\hline
1.3 & 금지 변수 사용 여부 확인 & \checkmark / \times & 신용정보법 \\
\hline
1.4 & 프록시 차별 리스크 평가 & \checkmark / \times & ECOA, 금융위 Guide라인 \\
\hline
\multicolumn{4}{|c|}{\textbf{2. 공정성 평가}} \\
\hline
2.1 & Adverse Impact 분석 (4/5 규칙) & \checkmark / \times & ECOA, 금융위 Guide라인 \\
\hline
2.2 & 집단별 성능 차이 분석 & \checkmark / \times & 금융위 Guide라인 \\
\hline
2.3 & Calibration 분석 & \checkmark / \times & SR 11-7 \\
\hline
2.4 & 편향 완화 조치 적용 (필요 시) & \checkmark / \times & 금융위 Guide라인 \\
\hline
\multicolumn{4}{|c|}{\textbf{3. 설명가능성}} \\
\hline
3.1 & 개별 결정 설명 생성 가능 & \checkmark / \times & 신용정보법 36조의2 \\
\hline
3.2 & Adverse Action Notice 템플릿 준비 & \checkmark / \times & ECOA Regulation B \\
\hline
3.3 & 고객 문의 Response 절차 수립 & \checkmark / \times & 금융위 Guide라인 \\
\hline
\multicolumn{4}{|c|}{\textbf{4. 내부 Control}} \\
\hline
4.1 & 모델 검증 (독립적 2차 검증) & \checkmark / \times & SR 11-7 \\
\hline
4.2 & 모델 인벤토리 등록 & \checkmark / \times & SR 11-7 \\
\hline
4.3 & 인간 감독 체계 구축 & \checkmark / \times & 시행령 14 \\
\hline
4.4 & 모니터링 체계 구축 & \checkmark / \times & SR 11-7 \\
\hline
\multicolumn{4}{|c|}{\textbf{5. 문서화}} \\
\hline
5.1 & 모델 카드 작성 & \checkmark / \times & 시행령 15 \\
\hline
5.2 & 공정성 평가 Report 작성 & \checkmark / \times & 시행령 15 \\
\hline
5.3 & 5년 보관 체계 확립 & \checkmark / \times & 시행령 15 \\
\hline
\end{tabularx}
\end{table}

\subsection{생성형 AI의 금융 부문 도입과 거버넌스}
\label{sec:5.1.5}

2024--2025년을 기점으로 대형 언어 모델(LLM)과 생성형 AI가 금융 산업 전반에
빠르게 확산되면서, 기존 AI 거버넌스 프레임워크로는 대응하기 어려운
새로운 유형의 리스크가 부상하고 있다.
금융위원회는 2025년 ``금융권 생성형 AI 활용 지원방안''을 발표하고,
기존 ``금융분야 AI 가이드라인''을 7대 원칙(거버넌스, 합법성, 보조수단성,
신뢰성, 금융안정성, 신의성실, 보안성)으로 개정하였다\cite{fsc_genai_2025}.
미국 FINRA 역시 2026년 연례 규제 감독 보고서(Annual Regulatory Oversight Report)에서
생성형 AI를 독립 섹션으로 다루며, 환각(Hallucination), 편향,
딥페이크 사기 등에 대한 감독 요구사항을 구체화하였다\cite{finra_2026_report}.

\paragraph{금융 LLM의 주요 리스크 유형}

\begin{table}[htbp]
\centering
\caption{금융 부문 생성형 AI 주요 리스크와 대응 방안}
\label{tab:genai_finance_risks}
\begin{tabularx}{\textwidth}{p{2.5cm}p{3.5cm}X}
\toprule
\textbf{리스크 유형} & \textbf{구체적 시나리오} & \textbf{거버넌스 대응 방안} \\
\midrule
환각(Hallucination) &
규정, 상품 조건, 금리 등에 대한 허위 정보 생성 &
RAG(Retrieval-Augmented Generation) 기반 응답 검증; 출처 명시 의무화; Human-in-the-loop 검토 \\
\midrule
딥페이크 사기 &
AI 생성 음성으로 콜센터 본인인증 우회; AI 합성 신분증으로 계좌 개설 &
다중 인증(MFA) 강화; 생체인증 라이브니스(Liveness) 검사; AI 생성 콘텐츠 탐지 모델 도입 \\
\midrule
시장 조작 &
LLM 기반 자동 매매에서의 담합적 행동; AI 생성 허위 뉴스를 통한 주가 조작 &
AI 거래 행동 감시 시스템; 이상 거래 탐지 알고리즘 고도화; AI 에이전트 행위 감사 추적(Audit Trail) \\
\midrule
개인정보 유출 &
LLM 프롬프트에 고객 개인정보 입력; 학습 데이터에서 민감 정보 추출 &
프롬프트 필터링; DLP(Data Loss Prevention) 연동; 내부망 전용 모델 운영 \\
\midrule
편향 증폭 &
LLM이 특정 고객 집단에 불리한 상품 추천 또는 상담 응대 차별 &
응답 공정성 정기 감사; 민감 속성별 응답 품질 비교 분석 \\
\bottomrule
\end{tabularx}
\end{table}


\paragraph{금융 생성형 AI 거버넌스 필수 요건}

FINRA의 2026년 보고서와 금융위원회의 개정 가이드라인을 종합하면,
금융 기관이 생성형 AI를 도입할 때 준수해야 할 핵심 요건은 다음과 같다:

\begin{enumerate}
 \item \textbf{사전 승인 체계}: 모든 생성형 AI 활용 사례(Use Case)는
 서면으로 목적, 데이터 소스, 모델/제공업체 선정, 통제 설계를 명시하여
 사전 승인을 받아야 한다.

 \item \textbf{고객 대면 출력의 인간 검증}: 고객에게 전달되거나
 의사결정에 영향을 미치는 모든 LLM 출력에 대해 Human-in-the-loop 검증을
 수행하고, 서명된 승인 기록을 보관해야 한다.

 \item \textbf{AI 에이전트 통제}: 거래를 실행하거나 행위를 수행할 수 있는
 AI 에이전트에 대해서는 좁은 범위의 권한 설정, 행위 감사 추적(Audit Trail),
 실행 전 명시적 인간 승인 체크포인트를 구현해야 한다.

 \item \textbf{환각 방지 설계}: 금융 규정, 상품 정보, 시장 데이터 등
 사실 기반 응답이 필요한 영역에서는 RAG 아키텍처를 적용하고,
 환각 발생률을 정량적으로 모니터링해야 한다.

 \item \textbf{기록 보존}: 생성형 AI와의 모든 상호작용(프롬프트, 응답)을
 감독 규정에 따라 기록하고 보존해야 한다.
\end{enumerate}


%----------------------------------------------------------
\section{의료.헬스케어 산업}
\label{sec:5.2}

\subsection{의료 AI의 규제 환경}
\label{sec:5.2.1}

의료 AI는 환자의 생명과 건강에 직접적인 영향을 미치므로, 
안전성 유효성에 대한 엄격한 규제를 받는다. 
동시에 의료 데이터의 민감성으로 인해 Privacy 보호 요구도 높다.

\begin{table}[htbp]
\centering
\caption{의료 AI 관련 주요 규제}
\label{tab:healthcare_regulations}
\begin{tabularx}{\textwidth}{p{3.5cm}p{2.5cm}X}
\toprule
\textbf{규제/Guide라인} & \textbf{관할} & \textbf{AI 관련 핵심 내용} \\
\midrule
의료기기법 & 
한국 & 
AI based 의료기기 허가/인증; SaMD(Software as a Medical Device) 규제 \\
\midrule
빅데이터 및 AI based 의료기기 허가심사 Guide라인 (식약처, 2024) & 
한국 & 
임상적 유효성, 알고리즘 검증, 학습 데이터 요구사항 \\
\midrule
개인정보 보호법 & 
한국 & 
민감정보(건강정보) 처리 제한; 가명처리 특례 \\
\midrule
생명윤리 및 안전에 관한 법률 & 
한국 & 
인체유래물, 유전정보 연구 규제; IRB 승인 \\
\midrule
AI 기본법 시행령 & 
한국 & 
의료를 고위험 AI로 분류 (제10조) \\
\midrule
EU AI Act & 
EU & 
의료기기를 고위험 AI로 분류; MDR/IVDR 연계 \\
\midrule
FDA AI/ML-Based SaMD Guidance & 
미국 & 
Good Machine Learning Practice (GMLP), 지속적 학습 모델 규제 \\
\midrule
HIPAA & 
미국 & 
Protected Health Information (PHI) 보호 \\
\bottomrule
\end{tabularx}
\end{table}


\paragraph{의료기기로서의 AI 분류}

한국 식품의약품안전처는 AI based 소프트웨어를 
의료기기로 분류하는 기준을 제시하고 있다:

\begin{itemize}
 \item \textbf{의료기기 해당}: 
 질병 진단, 예후 예측, 치료 계획 수립, 의료 영상 분석 등 
 의료 목적의 의사결정을 Support하는 AI
 
 \item \textbf{의료기기 비해당}: 
 행정 업무 자동화, 일정 관리, 일반 건강 정보 제공 등 
 의료 의사결정에 직접 관여하지 않는 AI
\end{itemize}


\subsection{의료 AI 주요 활용 사례와 특수 고려사항}
\label{sec:5.2.2}

\begin{table}[htbp]
\centering
\caption{의료 AI 활용 사례별 거버넌스 고려사항}
\label{tab:healthcare_ai_cases}
\begin{tabularx}{\textwidth}{p{2.5cm}p{2.5cm}p{2.5cm}X}
\toprule
\textbf{활용 Case} & \textbf{리스크 등급} & \textbf{주요 리스크} & \textbf{거버넌스 중점 사항} \\
\midrule
의료 영상 진단 (X-ray, CT, MRI) & 
고위험 (Class III) & 
오진, 위음성으로 인한 치료 지연 & 
임상 검증, 민감도/특이도 분석, 하위 집단 성능, 방사선 전문의 최종 판독 \\
\midrule
병리 조직 분석 & 
고위험 (Class III) & 
암 오진, 과잉/과소 치료 & 
다기관 검증, 병리 전문의 검토 필수, 불확실성 표시 \\
\midrule
임상 의사결정 Support (CDSS) & 
중-고위험 & 
부적절 권고, 알림 피로 & 
의료진 최종 결정권, 권고 근거 설명, 오버라이드 추적 \\
\midrule
환자 리스크 예측 (패혈증, 재입원 등) & 
중위험 & 
FN으로 인한 조기 개입 실패, FP로 인한 Resource 낭비 & 
임계값 Setup 근거, 집단별 성능, 조기 경보 프로토콜 \\
\midrule
신약 개발 (분자 설계, 후보 물질 탐색) & 
연구 단계 & 
비효과적 후보, 안전성 문제 & 
실험적 검증 필수, 해석가능성, 안전성 예측 \\
\midrule
정신건강 평가 & 
고위험 & 
자살 리스크 예측 실패, 낙인 효과 & 
극도의 민감성, 인간 전문가 개입 필수, Privacy 강화 \\
\bottomrule
\end{tabularx}
\end{table}


\paragraph{의료 AI의 특수한 윤리적 고려사항}

의료 AI는 일반적인 AI 윤리 원칙 외에 추가적인 고려가 필요하다:

\begin{enumerate}
 \item \textbf{환자 안전 우선(Primum Non Nocere)}: 
 정확도보다 안전이 우선; 불확실할 때는 보수적 결정
 
 \item \textbf{임상적 유효성(Clinical Validity)}: 
 기술적 성능만으로 불충분; 실제 임상 환경에서의 효과 입증 필요
 
 \item \textbf{건강 형평성(Health Equity)}: 
 인종, 성별, 사회경제적 지위에 따른 건강 격차 확대 방지; 
 학습 데이터의 대표성 확보
 
 \item \textbf{의료진-AI 협력}: 
 AI를 대체가 아닌 보조 Tool로 포지셔닝; 
 의료진의 전문성 AI의 장점 결합
 
 \item \textbf{환자 자율성(Patient Autonomy)}: 
 AI 활용에 대한 충분한 설명과 동의; 
 AI based 결정 거부 권리 보장
\end{enumerate}


\subsection{의료 AI 검증 프레임워크}
\label{sec:5.2.3}

\begin{lstlisting}[language=Python, caption={Healthcare AI Validation Framework}, label={lst:healthcare_validation}]
from dataclasses import dataclass, field
from typing import Dict, List, Optional, Tuple
import numpy as np
from sklearn.metrics import (
 roc_auc_score, precision_recall_curve, 
 confusion_matrix, accuracy_score
)

@dataclass
class ClinicalValidationConfig:
 """Healthcare AI Clinical Setup"""
 
 # Analysis Target
 subgroups: List[str] = field(default_factory=lambda: [
 'sex', # 
 'age_group', # 
 'ethnicity', # / 
 'comorbidity', # 'device_type', # ( AI)
 'institution', # Institution ( Institution )
 ])
 
 # Performance 
 min_sensitivity: float = 0.90 # Minimum 
 min_specificity: float = 0.80 # Minimum 
 max_subgroup_gap: float = 0.10 # Performance 
 
 # Interval
 confidence_level: float = 0.95


class MedicalAIValidator:
 """
 Healthcare AI Tool
 
 Guide , FDA GMLP Compliance Clinical Assessment 
 """
 
 def __init__(self, config: Optional[ClinicalValidationConfig] = None):
 self.config = config or ClinicalValidationConfig()
 
 def evaluate_diagnostic_performance(
 self,
 y_true: np.ndarray,
 y_prob: np.ndarray,
 operating_point: Optional[float] = None,
 ) -> Dict:
 """
 AI Performance Assessment
 
 , , PPV, NPV 95% Interval Calculation
 """
 
 # AUC
 auc = roc_auc_score(y_true, y_prob)
 
 # operating point (Youden's J)
 if operating_point is None:
 precision, recall, thresholds = precision_recall_curve(y_true, y_prob)
 # Sensitivity-Specificity based fpr_tpr_thresholds = self._calculate_roc_points(y_true, y_prob)
 operating_point = self._find_optimal_threshold(fpr_tpr_thresholds)
 
 # y_pred = (y_prob >= operating_point).astype(int)
 tn, fp, fn, tp = confusion_matrix(y_true, y_pred).ravel()
 
 # Major 
 sensitivity = tp / (tp + fn) if (tp + fn) > 0 else 0
 specificity = tn / (tn + fp) if (tn + fp) > 0 else 0
 ppv = tp / (tp + fp) if (tp + fp) > 0 else 0
 npv = tn / (tn + fn) if (tn + fn) > 0 else 0
 
 # 95% Interval (Wilson score interval)
 sensitivity_ci = self._wilson_ci(tp, tp + fn)
 specificity_ci = self._wilson_ci(tn, tn + fp)
 ppv_ci = self._wilson_ci(tp, tp + fp)
 npv_ci = self._wilson_ci(tn, tn + fn)
 
 return {
 'auc': auc,
 'operating_point': operating_point,
 'sensitivity': {'value': sensitivity, 'ci_95': sensitivity_ci},
 'specificity': {'value': specificity, 'ci_95': specificity_ci},
 'ppv': {'value': ppv, 'ci_95': ppv_ci},
 'npv': {'value': npv, 'ci_95': npv_ci},
 'confusion_matrix': {'tn': tn, 'fp': fp, 'fn': fn, 'tp': tp},
 'meets_sensitivity_threshold': sensitivity >= self.config.min_sensitivity,
 'meets_specificity_threshold': specificity >= self.config.min_specificity,
 }
 
 def evaluate_subgroup_performance(
 self,
 y_true: np.ndarray,
 y_prob: np.ndarray,
 subgroup_data: Dict[str, np.ndarray],
 operating_point: float,
 ) -> Dict:
 """
 Performance Analysis
 
 Equity Assessment: , , Performance Analysis
 """
 
 results = {}
 
 for subgroup_name, subgroup_labels in subgroup_data.items():
 unique_groups = np.unique(subgroup_labels)
 group_performance = {}
 
 for group in unique_groups:
 mask = subgroup_labels == group
 if mask.sum() < 30: # Minimum Number of samples
 continue
 
 group_y_true = y_true[mask]
 group_y_prob = y_prob[mask]
 
 perf = self.evaluate_diagnostic_performance(
 group_y_true, group_y_prob, operating_point
 )
 
 group_performance[str(group)] = {
 'n': mask.sum(),
 'sensitivity': perf['sensitivity']['value'],
 'specificity': perf['specificity']['value'],
 'auc': perf['auc'],
 }
 
 # Performance Analysis
 if len(group_performance) >= 2:
 sensitivities = [g['sensitivity'] for g in group_performance.values()]
 specificities = [g['specificity'] for g in group_performance.values()]
 
 sens_gap = max(sensitivities) - min(sensitivities)
 spec_gap = max(specificities) - min(specificities)
 
 results[subgroup_name] = {
 'groups': group_performance,
 'sensitivity_gap': sens_gap,
 'specificity_gap': spec_gap,
 'equity_concern': sens_gap > self.config.max_subgroup_gap or spec_gap > self.config.max_subgroup_gap,
 }
 
 return results
 
 def _wilson_ci(
 self,
 successes: int,
 trials: int,
 z: float = 1.96,
 ) -> Tuple[float, float]:
 """Wilson score confidence interval"""
 
 if trials == 0:
 return (0.0, 0.0)
 
 p = successes / trials
 denominator = 1 + z2 / trials
 center = (p + z2 / (2 * trials)) / denominator
 margin = z * np.sqrt(p * (1 - p) / trials + z**2 / (4 * trials2)) / denominator
 
 return (max(0, center - margin), min(1, center + margin))
 
 def _calculate_roc_points(
 self,
 y_true: np.ndarray,
 y_prob: np.ndarray,
 ) -> List[Tuple[float, float, float]]:
 """ROC Calculation"""
 
 thresholds = np.unique(y_prob)
 points = []
 
 for thresh in thresholds:
 y_pred = (y_prob >= thresh).astype(int)
 tn, fp, fn, tp = confusion_matrix(y_true, y_pred).ravel()
 
 tpr = tp / (tp + fn) if (tp + fn) > 0 else 0
 fpr = fp / (fp + tn) if (fp + tn) > 0 else 0
 
 points.append((fpr, tpr, thresh))
 
 return points
 
 def _find_optimal_threshold(
 self,
 roc_points: List[Tuple[float, float, float]],
 ) -> float:
 """Youden's J statistic based """
 
 best_j = -1
 best_thresh = 0.5
 
 for fpr, tpr, thresh in roc_points:
 j = tpr - fpr # Youden's J
 if j > best_j:
 best_j = j
 best_thresh = thresh
 
 return best_thresh
 
 def generate_clinical_validation_report(
 self,
 model_name: str,
 indication: str,
 performance: Dict,
 subgroup_analysis: Dict,
 ) -> str:
 """ Clinical Report Generation"""
 
 report = f"""
# Healthcare AI Clinical ReportModel : {model_name} : {indication} : Data AI based Healthcare Guide 

---

## 1. Performance

| | | 95% CI | |
|--------|-----|--------|----------|
| AUC | {performance['auc']:.4f} | - | - |
| | {performance['sensitivity']['value']:.4f} | [{performance['sensitivity']['ci_95'][0]:.4f}, {performance['sensitivity']['ci_95'][1]:.4f}] | {'[PASS]' if performance['meets_sensitivity_threshold'] else '[FAIL]'} (>={self.config.min_sensitivity}) |
| | {performance['specificity']['value']:.4f} | [{performance['specificity']['ci_95'][0]:.4f}, {performance['specificity']['ci_95'][1]:.4f}] | {'[PASS]' if performance['meets_specificity_threshold'] else '[FAIL]'} (>={self.config.min_specificity}) |
| (PPV) | {performance['ppv']['value']:.4f} | [{performance['ppv']['ci_95'][0]:.4f}, {performance['ppv']['ci_95'][1]:.4f}] | - |
| (NPV) | {performance['npv']['value']:.4f} | [{performance['npv']['ci_95'][0]:.4f}, {performance['npv']['ci_95'][1]:.4f}] | - |

### | | | |
|--|----------|----------|
| | {performance['confusion_matrix']['tp']} (TP) | {performance['confusion_matrix']['fn']} (FN) |
| | {performance['confusion_matrix']['fp']} (FP) | {performance['confusion_matrix']['tn']} (TN) |

---

## 2. Analysis ( Equity)

"""
 
 equity_concerns = []
 
 for subgroup, result in subgroup_analysis.items():
 report += f"### {subgroup}\n\n"
 
 if result.get('equity_concern', False):
 equity_concerns.append(subgroup)
 report += f"[WARN]Equity : Performance ({self.config.max_subgroup_gap:.0%}) \n\n"
 
 report += "| | N | | | AUC |\n"
 report += "|------|---|--------|--------|-----|\n"
 
 for group, stats in result.get('groups', {}).items():
 report += f"| {group} | {stats['n']} | {stats['sensitivity']:.4f} | {stats['specificity']:.4f} | {stats['auc']:.4f} |\n"
 
 report += f"\n- : {result['sensitivity_gap']:.4f}\n"
 report += f"- : {result['specificity_gap']:.4f}\n\n"
 
 report += """
---

## 3. Assessment """
 
 if not performance['meets_sensitivity_threshold']:
 report += "[FAIL] : Impact \n\n"
 
 if not performance['meets_specificity_threshold']:
 report += "[WARN] : Addition \n\n"
 
 if equity_concerns:
 report += f"[WARN] Equity : {', '.join(equity_concerns)} Performance \n\n"
 
 if (performance['meets_sensitivity_threshold'] and 
 performance['meets_specificity_threshold'] and 
 not equity_concerns):
 report += "[PASS] Clinical **\n\n"
 
 report += """
---

## 4. 1. AI Healthcare Tool , Final Healthcare .
2. Learning Data , Performance .
3. Clinical AI days , Healthcare Decision .

---

* Report 15 5years .*
"""
 
 return report
\end{lstlisting}


\subsection{의료 AI 거버넌스 체크리스트}
\label{sec:5.2.4}

\begin{table}[htbp]
\centering
\caption{[실무 도구 5-2] 의료 AI 거버넌스 체크리스트}
\label{tab:healthcare_checklist}
\begin{tabularx}{\textwidth}{|p{0.5cm}|X|p{1.5cm}|p{3cm}|}
\hline
\textbf{No.} & \textbf{Check 항목} & \textbf{결과} & \textbf{근거 규정} \\
\hline
\multicolumn{4}{|c|}{\textbf{1. 규제 분류}} \\
\hline
1.1 & 의료기기 해당 여부 결정 & \checkmark / \times & 의료기기법 \\
\hline
1.2 & 리스크 등급 분류 (Class I/II/III) & \checkmark / \times & 의료기기법 \\
\hline
1.3 & 허가/인증 경로 결정 & \checkmark / \times & 식약처 Guide라인 \\
\hline
\multicolumn{4}{|c|}{\textbf{2. 데이터 및 Learning}} \\
\hline
2.1 & IRB 승인 (해당 시) & \checkmark / \times & 생명윤리법 \\
\hline
2.2 & 학습 데이터 대표성 확보 & \checkmark / \times & 식약처 Guide라인 \\
\hline
2.3 & 데이터 품질 관리 & \checkmark / \times & GMLP \\
\hline
2.4 & 개인정보 비식별화/가명처리 & \checkmark / \times & 개인정보보호법 \\
\hline
\multicolumn{4}{|c|}{\textbf{3. 임상 검증}} \\
\hline
3.1 & 분석적 검증 수행 & \checkmark / \times & 식약처 Guide라인 \\
\hline
3.2 & 임상적 검증 수행 & \checkmark / \times & 식약처 Guide라인 \\
\hline
3.3 & 하위 집단 성능 분석 & \checkmark / \times & 건강 형평성 \\
\hline
3.4 & 다기관 검증 (해당 시) & \checkmark / \times & 식약처 Guide라인 \\
\hline
\multicolumn{4}{|c|}{\textbf{4. 안전 및 유효성}} \\
\hline
4.1 & 민감도/특이도 기준 충족 & \checkmark / \times & 식약처 Guide라인 \\
\hline
4.2 & 사용 목적 명확화 & \checkmark / \times & 의료기기법 \\
\hline
4.3 & 의료진 최종 결정 보장 & \checkmark / \times & 시행령 14 \\
\hline
4.4 & 불확실성 표시 기능 & \checkmark / \times & GMLP \\
\hline
\multicolumn{4}{|c|}{\textbf{5. 시판 후 관리}} \\
\hline
5.1 & 성능 모니터링 계획 & \checkmark / \times & 의료기기법 \\
\hline
5.2 & 이상사례 보고 체계 & \checkmark / \times & 의료기기법 \\
\hline
5.3 & 재학습/업데이트 절차 & \checkmark / \times & FDA GMLP \\
\hline
\end{tabularx}
\end{table}

%----------------------------------------------------------
\section{Public 부문 (Public Sector)}
\label{sec:5.3}

\subsection{Public AI의 규제 환경 및 Principle}
\label{sec:5.3.1}

공공 부문에서의 AI 활용은 시민의 권리와 의무에 직접적인 영향을 미치며, 세금으로 운영된다는 점에서 투명성, 공정성, 책임성이 극도로 강조된다. 전 세계적으로 공공 부문의 AI 활용을 규율하는 Guide라인이 빠르게 정립되고 있다.

\begin{table}[htbp]
\centering
\caption{Public 부문 AI 관련 주요 규제 및 정책}
\label{tab:public_regulations}
\begin{tabularx}{\textwidth}{p{3.5cm}p{2.5cm}X}
\toprule
\textbf{규제/정책} & \textbf{국가/Institution} & \textbf{핵심 내용} \\
\midrule
AI in Government Act & 미국 & 연방 기관의 AI 활용 기준 마련; 역량 강화 및 거버넌스 프레임워크 \\
\midrule
Executive Order 14110 & 미국 & 안전하고 신뢰할 수 있는 AI 개발 및 사용; 권리 보호 및 형평성 강화 \\
\midrule
PublicInstitution AI 도입.활용 Guide라인 & 한국 & 공공성, 투명성, 책임성 3대 원칙; 리스크 평가 및 모니터링 절차 \\
\midrule
지능정보화 기본법 & 한국 & 공공기관의 지능정보사회 구현 의무; 정보격차 해소 및 안전성 확보 \\
\midrule
EU AI Act (Public Sector) & EU & 공공 서비스에서의 생체 인식 금지; 고위험 AI에 대한 엄격한 투명성 의무 \\
\bottomrule
\end{tabularx}
\end{table}

\paragraph{Public AI의 핵심 가치}
공공 부문 AI 거버넌스는 다음과 같은 민주적 가치를 수호해야 한다:
\begin{itemize}
 \item \textbf{적법 절차(Due 프로세스)}: AI based 결정이 시민의 권리를 침해할 경우, 이에 대한 설명과 불복 절차가 보장되어야 한다.
 \item \textbf{비차별 및 형평성}: 특정 계층이나 집단이 서비스에서 소외되거나 부당한 대우를 받지 않도록 알고리즘 Fairness을 확보해야 한다.
 \item \textbf{투명성(투명성)}: 어떤 데이터가 사용되었고, 어떤 논리로 결과가 도출되었는지 시민이 이해할 수 있는 수준으로 공개해야 한다.
\end{itemize}


\subsection{Public 부문 주요 활용 사례와 리스크}
\label{sec:5.3.2}

\begin{table}[htbp]
\centering
\caption{Public AI 활용 사례별 거버넌스 고려사항}
\label{tab:public_ai_cases}
\begin{tabularx}{\textwidth}{p{2.5cm}p{2.5cm}p{2.5cm}X}
\toprule
\textbf{활용 Case} & \textbf{리스크 등급} & \textbf{주요 리스크} & \textbf{거버넌스 중점 사항} \\
\midrule
복지 혜택 대상자 선정 & 고위험 & 오분류로 인한 수급 탈락, 차별 & 알고리즘 공정성 검증, 이의 신청 절차, 인간의 최종 승인 \\
\midrule
치안 및 범죄 예측 (Predictive Policing) & 매우 높음 & 편향된 데이터로 인한 특정 지역/인종 감시 강화 & 투명성 Report 발간, 외부 감시 체계, 경찰권 남용 방지 \\
\midrule
교통 및 도시 관리 (Smart City) & 중위험 & Privacy 침해, 시스템 오작동 & 개인정보 비식별화, 보안 사고 Response 계획, 시민 참여형 거버넌스 \\
\midrule
민원 자동 응답 (Chatbot) & 저-중위험 & 잘못된 정보 제공, 편향적 발언 & 답변 품질 모니터링, Guide라인 준수, 답변 근거 제시 \\
\bottomrule
\end{tabularx}
\end{table}


\subsection{실무 도구: Public AI 투명성 Report 생성 도구}
\label{sec:5.3.3}

공공 부문에서는 AI 시스템의 운영 현황을 시민에게 공개하는 것이 중요하다. 다음은 모델의 투명성을 확보하기 위한 Report 자동 생성 Tool의 예시이다.

\begin{lstlisting}[language=Python, caption={Public AI Transparency Report Generator}, label={lst:public_transparency}]
import json
from datetime import datetime
from typing import Dict, List, Any

class PublicAITransparencyReport:
 """
 PublicInstitution AI System Report Generator
 
 EO 14110 PublicInstitution AI Guide Compliance Support
 """
 
 def __init__(self, system_name: str, department: str):
 self.system_info = {
 "system_name": system_name,
 "department": department,
 "generation_date": datetime.now().isoformat(),
 "data_sources": [],
 "algorithms": [],
 "impact_assessment": {},
 "human_oversight": ""
 }

 def add_data_source(self, name: str, source: str, collection_method: str):
 self.system_info["data_sources"].append({
 "name": name,
 "source": source,
 "method": collection_method
 })

 def set_algorithm_details(self, model_type: str, framework: str, version: str):
 self.system_info["algorithms"].append({
 "type": model_type,
 "framework": framework,
 "version": version
 })

 def set_impact_assessment(self, target_population: str, potential_risks: List[str], mitigations: List[str]):
 self.system_info["impact_assessment"] = {
 "target": target_population,
 "risks": potential_risks,
 "mitigations": mitigations
 }

 def set_human_oversight(self, role: str, process: str):
 self.system_info["human_oversight"] = f" : {role}, : {process}"

 def generate_report(self) -> str:
 report = f"# AI System Report: {self.system_info['system_name']}\n\n"
 report += f" : {self.system_info['department']}\n"
 report += f"days: {self.system_info['generation_date']}\n\n"
 
 report += "## 1. Data Usage \n"
 for ds in self.system_info["data_sources"]:
 report += f"- {ds['name']} ({ds['source']}): {ds['method']}\n"
 
 report += "\n## 2. information\n"
 for alg in self.system_info["algorithms"]:
 report += f"- {alg['type']} (ver {alg['version']}), Framework: {alg['framework']}\n"
 
 report += "\n## 3. Impact Assessment \n"
 impact = self.system_info["impact_assessment"]
 report += f"- Target : {impact.get('target', 'N/A')}\n"
 report += "- Major Risk:\n"
 for risk in impact.get('risks', []):
 report += f" - {risk}\n"
 report += "- Response :\n"
 for mit in impact.get('mitigations', []):
 report += f" - {mit}\n"
 
 report += f"\n## 4. \n{self.system_info['human_oversight']}\n"
 
 return report

# Example
reporter = PublicAITransparencyReport(" AI", " Support ")
reporter.add_data_source(" / information", " / ", " Integration")
reporter.set_algorithm_details("Random Forest Regressor", "Scikit-learn", "1.2.0")
reporter.set_impact_assessment(
 " ",
 [" ", " "],
 [" Weights ", " "]
)
reporter.set_human_oversight(" ", "AI list based Final Target ")

print(reporter.generate_report())
\end{lstlisting}


\subsection{Public 부문 AI 거버넌스 체크리스트}
\label{sec:5.3.4}

\begin{table}[htbp]
\centering
\caption{[실무 도구 5-3] 공공 부문 AI 거버넌스 체크리스트}
\label{tab:public_checklist}
\begin{tabularx}{\textwidth}{|p{0.5cm}|X|p{1.5cm}|p{3cm}|}
\hline
\textbf{No.} & \textbf{Check 항목} & \textbf{결과} & \textbf{근거 규정} \\
\hline
\multicolumn{4}{|c|}{\textbf{1. 법적 근거 및 절차}} \\
\hline
1.1 & AI 도입의 법적 근거 확보 (법령/조례) & \checkmark / \times & 지능정보화 기본법 \\
\hline
1.2 & 행정 절차법상 적법 절차 준수 여부 & \checkmark / \times & 행정절차법 \\
\hline
\multicolumn{4}{|c|}{\textbf{2. 투명성 및 설명가능성}} \\
\hline
2.1 & 알고리즘 개요 및 사용 데이터 공개 & \checkmark / \times & 지능정보화 기본법 62 \\
\hline
2.2 & 결과 도출의 주요 요인 설명 가능 여부 & \checkmark / \times & 시행령 14 \\
\hline
2.3 & 투명성 Report 정기 발간 계획 & \checkmark / \times & Guide라인 \\
\hline
\multicolumn{4}{|c|}{\textbf{3. 공정성 및 차별 방지}} \\
\hline
3.1 & 데이터 수집 과정 편향 Check & \checkmark / \times & Guide라인 \\
\hline
3.2 & 취약 계층에 대한 알고리즘 공정성 평가 & \checkmark / \times & 시행령 13 \\
\hline
3.3 & 정보 취약 계층의 접근성 보장 & \checkmark / \times & 지능정보화 기본법 46 \\
\hline
\multicolumn{4}{|c|}{\textbf{4. 책임성 및 인간 제어}} \\
\hline
4.1 & 최종 결정권자의 인간 지정 (Human-in-the-loop) & \checkmark / \times & 시행령 14 \\
\hline
4.2 & AI 오류 발생 시 시민 피해보상 체계 & \checkmark / \times & 국가배상법 \\
\hline
4.3 & 이의 신청 및 불복 절차 마련 & \checkmark / \times & 행정절차법 \\
\hline
\multicolumn{4}{|c|}{\textbf{5. 정보 보안 및 개인정보}} \\
\hline
5.1 & Public Data Security Guide라인 준수 & \checkmark / \times & 국가 정보보안 기본지침 \\
\hline
5.2 & 민감 행정 데이터의 비식별화 처리 & \checkmark / \times & 개인정보 보호법 \\
\hline
\end{tabularx}
\end{table}

\subsection{공공 부문 생성형 AI 도입과 거버넌스}
\label{sec:5.3.5}

디지털플랫폼정부위원회는 2025년 ``공공부문 초거대 AI 도입.활용 가이드라인 2.0''을
발표하며, 중앙부처, 지방자치단체, 공공기관의 생성형 AI 활용에 대한
체계적 지침을 제시하였다\cite{dpg_guideline_2025}.
공공 부문의 생성형 AI 도입은 대국민 서비스 혁신과 행정 효율화라는
기회를 제공하지만, 동시에 행정 신뢰성, 개인정보 보호, 정보 접근성 등
고유한 거버넌스 과제를 수반한다.

\paragraph{시민 대면 AI 챗봇의 거버넌스 요건}

공공기관의 민원 응대 챗봇은 시민의 권리.의무에 영향을 미치는 정보를
제공한다는 점에서, 일반적인 고객 서비스 챗봇보다 높은 수준의
거버넌스가 요구된다.

\begin{table}[htbp]
\centering
\caption{공공 민원 챗봇 거버넌스 요건}
\label{tab:public_chatbot_governance}
\begin{tabularx}{\textwidth}{p{2.5cm}p{3cm}X}
\toprule
\textbf{거버넌스 영역} & \textbf{요건} & \textbf{구현 방안} \\
\midrule
정확성 보장 &
법령, 제도 안내의 정확성 확보 &
RAG 기반 공식 법령 DB 연동; 정기적 정확도 평가(월 1회 이상); 오답률 임계값(예: 2\% 이하) 설정 및 모니터링 \\
\midrule
책임성 확보 &
AI 응답에 대한 행정적 책임 소재 명확화 &
AI 응답 면책 고지(Disclaimer) 표시; ``이 응답은 AI가 생성한 참고 정보입니다'' 안내; 담당 공무원 연결 경로 제공 \\
\midrule
접근성 보장 &
정보 취약 계층의 서비스 이용 보장 &
음성 안내, 큰 글씨 모드, 다국어 지원; 고령자.장애인 전용 에스컬레이션 채널; 비AI 민원 채널 병행 운영 \\
\midrule
개인정보 보호 &
시민 상담 내역의 안전한 처리 &
대화 로그 비식별화 처리; 민감 정보(주민번호 등) 자동 마스킹; 보존 기간 설정 및 자동 삭제 \\
\midrule
투명성 &
AI 활용 사실의 시민 고지 &
챗봇 진입 시 AI 활용 사실 명시; 응답 생성 근거(출처 법령) 함께 제시; 연간 운영 현황 공개 \\
\bottomrule
\end{tabularx}
\end{table}


\paragraph{AI 생성 행정문서의 관리 기준}

생성형 AI를 활용한 행정문서 작성이 확산됨에 따라,
문서의 정확성, 법적 유효성, 책임 소재에 관한 새로운 기준이 필요하다.
개인정보보호위원회는 2025년 8월 ``생성형 AI 개발.활용을 위한
개인정보 처리 안내서''를 발표하여 공공기관의 AI 활용 시
개인정보 보호 기준을 제시하였다\cite{pipc_genai_2025}.

주요 관리 기준은 다음과 같다:

\begin{enumerate}
 \item \textbf{문서 유형별 AI 활용 범위 설정}:
 법적 구속력이 있는 처분서, 통지서 등은 AI 초안을 참고 자료로만 활용하고
 반드시 담당 공무원이 최종 검토.서명해야 한다.
 내부 보고서, 회의록 요약 등 비공식 문서는 AI 생성 후 검토 절차를
 간소화할 수 있다.

 \item \textbf{AI 생성 표시 의무}:
 AI가 초안을 작성하거나 내용 생성에 관여한 행정문서에는
 해당 사실을 메타데이터 또는 문서 하단에 명시해야 한다.

 \item \textbf{민감 데이터 입력 금지}:
 생성형 AI 프롬프트에 개인정보, 보안 등급 정보, 수사 관련 정보 등
 민감 행정 데이터를 입력하는 것을 원칙적으로 금지하며,
 불가피한 경우 내부망 전용 모델만 사용한다.

 \item \textbf{버전 관리 및 감사 추적}:
 AI 생성 초안과 최종 확정본의 변경 이력을 관리하고,
 감사 시 AI 관여 범위를 확인할 수 있도록 기록을 보존한다.
\end{enumerate}


%----------------------------------------------------------
\section{제조 및 산업 현장 (Manufacturing \& Industrial)}
\label{sec:5.4}

\subsection{제조 AI의 규제 환경 및 안전 표준}
\label{sec:5.4.1}

제조 산업에서의 AI 활용은 생산성 향상과 비용 절감을 주도하지만, 물리적 안전(안전성)과 직결된다는 점에서 독특한 규제 환경을 가진다. 특히 인간과 로봇이 협업하는 환경(Cobot)에서의 AI는 더욱 엄격한 안전 인증을 요구한다.

\begin{table}[htbp]
\centering
\caption{제조 AI 관련 주요 규제 및 안전 표준}
\label{tab:manufacturing_regulations}
\begin{tabularx}{\textwidth}{p{3.5cm}p{2.5cm}X}
\toprule
\textbf{규제/표준} & \textbf{관할/기구} & \textbf{핵심 내용} \\
\midrule
EU AI Act (Manufacturing) & EU & 기기 안전 부품으로 사용되는 AI를 고위험으로 분류; 적합성 평가 의무 \\
\midrule
ISO/IEC 23894 & ISO/IEC & AI 리스크 관리 표준; 제조 공정 불확실성 관리 가이드 제시 \\
\midrule
ISO 10218-1/2 & ISO & 산업용 로봇의 안전 요구사항; AI based 충돌 방지 시스템 인증 기준 \\
\midrule
중대재해 처벌법 & 한국 & AI 오작동으로 인한 인명 사고 시 경영책임자 처벌 가능성; 안전 보건 확보 의무 \\
\midrule
NIST AI RMF (Industrial Profile) & 미국 & 제조 현장의 특수성을 반영한 AI 리스크 관리 실무 지침 \\
\bottomrule
\end{tabularx}
\end{table}

\paragraph{제조 AI의 안전 최우선 Principle}
제조 현장의 AI 거버넌스는 'Fail-Safe' 설계를 핵심으로 한다:
\begin{itemize}
 \item \textbf{물리적 안전 확보}: AI의 결정 착오가 장비의 파손이나 작업자의 부상으로 이어지지 않도록 하드웨어 수준의 인터락(Interlock) 시스템과 연계해야 한다.
 \item \textbf{실시간성(Real-time) 보장}: 공정 제어 AI는 정해진 시간 내에 반응해야 하며, 지연(Latency)이 안전 사고로 이어지는 것을 방지해야 한다.
 \item \textbf{신뢰성 및 가용성}: 가혹한 산업 환경(고온, 소음, 진동)에서도 AI 모델이 일관된 성능을 유지하는지 검증해야 한다.
\end{itemize}


\subsection{제조 부문 주요 활용 사례와 리스크}
\label{sec:5.4.2}

\begin{table}[htbp]
\centering
\caption{제조 AI 활용 사례별 거버넌스 고려사항}
\label{tab:manufacturing_ai_cases}
\begin{tabularx}{\textwidth}{p{2.5cm}p{2.5cm}p{2.5cm}X}
\toprule
\textbf{활용 Case} & \textbf{리스크 등급} & \textbf{주요 리스크} & \textbf{거버넌스 중점 사항} \\
\midrule
예측 보전 (Predictive Maintenance) & 중위험 & 오진으로 인한 장기 가동 중단, 비용 발생 & 잔여 수명(RUL) 예측 오차 관리, 허위 경보(False Alarm) 최소화 \\
\midrule
자율 이동 로봇 (AMR/AGV) & 고위험 & 작업자 충돌, 물류 흐름 방해 & 센서 융합 데이터의 신뢰성, 비상 정지 기능 검증, 경로 계획 안전성 \\
\midrule
품질 검사 (Vision Inspection) & 저-중위험 & 불량품 출하, 과다 Retire & 검출 확률(PoD) 통계적 관리, 조명/환경 변화에 대한 강건성 \\
\midrule
공정 최적화 (Process 최적화) & 중위험 & 제어 루프 불안정, 비정상 상태 유발 & 운영 한계(Operational Envelope) Setup, 루트 코즈 분석(RCA) 기능 \\
\bottomrule
\end{tabularx}
\end{table}


\subsection{실무 도구: 제조 AI 안전성 평가 도구}
\label{sec:5.4.3}

제조 현장에서는 AI 모델의 성능보다 '안전 한계' 내에서의 작동 여부가 중요하다. 다음은 모델의 안전 임계값을 모니터링하고 사고를 예방하기 위한 평가 도구 예시이다.

\begin{lstlisting}[language=Python, caption={Manufacturing AI Safety Assessment Tool}, label={lst:manufacturing_safety}]
from dataclasses import dataclass
from typing import List, Tuple
import time

@dataclass
class ManufacturingAISafetyConfig:
 """Manufacturing AI Setup"""
 max_latency_ms: float = 50.0 # safety_buffer: float = 0.2 # (20%)
 reliability_threshold: float = 0.999 # emergency_stop_signals: List[str] = None

class ManufacturingAISafetyEvaluator:
 """
 Manufacturing AI Assessment 
 
 ISO 10218 EU AI Act Compliance Support
 """
 
 def __init__(self, config: ManufacturingAISafetyConfig):
 self.config = config
 self.safety_logs = []

 def check_operational_envelope(self, predicted_value: float, limits: Tuple[float, float]) -> bool:
 """
 AI Operation (Operational Envelope) """
 lower_limit, upper_limit = limits
 # Calculation
 safe_lower = lower_limit * (1 + self.config.safety_buffer) if lower_limit > 0 else lower_limit * (1 - self.config.safety_buffer)
 safe_upper = upper_limit * (1 - self.config.safety_buffer)
 
 is_safe = safe_lower <= predicted_value <= safe_upper
 
 if not is_safe:
 self._trigger_safety_alert(f" : {predicted_value} Scope ")
 
 return is_safe

 def measure_realtime_performance(self, start_time: float, end_time: float) -> bool:
 """
 AI """
 latency = (end_time - start_time) * 1000 # ms
 if latency > self.config.max_latency_ms:
 self._trigger_safety_alert(f" : {latency:.2f}ms ( : {self.config.max_latency_ms}ms)")
 return False
 return True

 def _trigger_safety_alert(self, message: str):
 log_entry = f"[{time.strftime('%Y-%m-%d %H:%M:%S')}] [SAFETY ALERT] {message}"
 self.safety_logs.append(log_entry)
 print(log_entry)
 # PLC(Logical Controller) # Example
config = ManufacturingAISafetyConfig(max_latency_ms=30.0, safety_buffer=0.15)
evaluator = ManufacturingAISafetyEvaluator(config)

# inference_start = time.time()
# ... AI Model ...
time.sleep(0.02) # 20ms inference_end = time.time()

# 1. is_realtime = evaluator.measure_realtime_performance(inference_start, inference_end)

# 2. Scope ( : )
torque_prediction = 85.5
operating_limit = (0.0, 100.0)
is_within_envelope = evaluator.check_operational_envelope(torque_prediction, operating_limit)

if not is_within_envelope or not is_realtime:
 print("System .")
\end{lstlisting}


\subsection{제조 부문 AI 거버넌스 체크리스트}
\label{sec:5.4.4}

\begin{table}[htbp]
\centering
\caption{[실무 도구 5-4] 제조 부문 AI 거버넌스 체크리스트}
\label{tab:manufacturing_checklist}
\begin{tabularx}{\textwidth}{|p{0.5cm}|X|p{1.5cm}|p{3cm}|}
\hline
\textbf{No.} & \textbf{Check 항목} & \textbf{결과} & \textbf{근거 규정} \\
\hline
\multicolumn{4}{|c|}{\textbf{1. 하드웨어 및 안전 인증}} \\
\hline
1.1 & AI 탑재 장비의 ISO 안전 인증 획득 여부 & \checkmark / \times & ISO 10218, ISO 13849 \\
\hline
1.2 & 독립적인 하드웨어 비상 정지(E-Stop) 연계 & \checkmark / \times & 기계 안전 표준 \\
\hline
1.3 & AI 소프트웨어 고장 시의 Fail-safe 모드 작동 & \checkmark / \times & EU AI Act 부속서 \\
\hline
\multicolumn{4}{|c|}{\textbf{2. 실시간성 및 강건성}} \\
\hline
2.1 & 추론 지연 시간(Latency)의 최악 상황 검증 & \checkmark / \times & Guide라인 \\
\hline
2.2 & 산업적 노이즈(Data 왜곡)에 대한 강건성 테스트 & \checkmark / \times & ISO 23894 \\
\hline
2.3 & 네트워크 단절 시 로컬 제어권 유지 로직 & \checkmark / \times & Guide라인 \\
\hline
\multicolumn{4}{|c|}{\textbf{3. 데이터 및 Security}} \\
\hline
3.1 & 현장 운영 Data(OT)의 보안 및 무결성 확보 & \checkmark / \times & IEC 62443 \\
\hline
3.2 & 센서 데이터의 이상치(Outlier) 감지 필터링 & \checkmark / \times & 품질 관리 표준 \\
\hline
3.3 & 모델 업데이트 전 현장 적응성 테스트(PoC) & \checkmark / \times & Guide라인 \\
\hline
\multicolumn{4}{|c|}{\textbf{4. 책임 및 인적 Resource}} \\
\hline
4.1 & AI 사고 시 책임 소재 규명 절차 (블랙박스) & \checkmark / \times & 중대재해처벌법 \\
\hline
4.2 & 현장 작업자 대상 AI 시스템 작동/위험 교육 & \checkmark / \times & 산업안전보건법 \\
\hline
4.3 & AI 권고 사항에 대한 감독자 승인 절차 & \checkmark / \times & 시행령 14 \\
\hline
\end{tabularx}
\end{table}

%----------------------------------------------------------
\section{소결: 산업별 거버넌스의 통합과 차별화}
\label{sec:5.5}

본 장에서 살펴본 금융, 의료, 공공, 제조의 네 가지 산업은 AI 거버넌스의 공통적인 토대(Data 보호, 투명성, 책임성) 위에서 각기 다른 최우선 가치를 추구하고 있다.

금융은 \textbf{공정성}과 \textbf{설명가능성}을 통해 신뢰를 구축하고, 의료는 \textbf{환자 안전}과 \textbf{임상적 유효성}을 절대적 기준으로 삼는다. 공공 부문은 \textbf{민주적 가치}와 \textbf{적법 절차}를 수호하며, 제조 산업은 \textbf{물리적 안전}과 \textbf{신뢰성}에 집중한다.

\subsection{산업 간 거버넌스 비교 분석}
\label{sec:5.5.1}

네 산업의 거버넌스 특성을 체계적으로 비교하면,
공통점과 차이점이 더욱 명확해진다.
표~\ref{tab:cross_industry_comparison}은
각 산업의 핵심 거버넌스 차원을 종합적으로 비교한 것이다.

\begin{table}[htbp]
\centering
\caption{산업별 AI 거버넌스 종합 비교}
\label{tab:cross_industry_comparison}
\begin{tabularx}{\textwidth}{p{2.2cm}p{2.8cm}p{2.8cm}p{2.8cm}p{2.8cm}}
\toprule
\textbf{거버넌스 차원} & \textbf{금융} & \textbf{의료} & \textbf{공공} & \textbf{제조} \\
\midrule
최우선 가치 &
공정성, 설명가능성 &
환자 안전, 임상 유효성 &
민주적 가치, 적법 절차 &
물리적 안전, 신뢰성 \\
\midrule
리스크 영향 범위 &
경제적 기회 (대출, 보험, 투자) &
생명 및 건강 &
시민 권리 및 공적 자원 배분 &
작업자 신체 안전 및 생산 연속성 \\
\midrule
주요 규제 체계 &
금융위 가이드라인, SR 11-7, ECOA &
식약처 가이드라인, 의료기기법, FDA GMLP &
지능정보화 기본법, 행정절차법 &
ISO 10218, 중대재해처벌법, IEC 62443 \\
\midrule
설명가능성 수준 &
개별 결정 수준 (Adverse Action Notice) &
임상적 근거 기반 설명 &
시민 이해 가능 수준 공개 &
고장 원인 분석(RCA) 수준 \\
\midrule
인간 감독 형태 &
모델 검증 위원회, 2차 검증 &
의료진 최종 판독/결정 &
담당 공무원 최종 승인 &
현장 감독자 승인, 비상 정지 \\
\midrule
생성형 AI 핵심 과제 &
환각 방지, 딥페이크 사기 대응 &
임상 정보 정확성, 환자 동의 &
행정문서 정확성, 시민 고지 &
공정 파라미터 생성의 안전성 \\
\midrule
모니터링 주기 &
실시간(FDS) + 분기별 공정성 감사 &
시판 후 지속 감시 + 이상사례 즉시 보고 &
연간 투명성 보고서 + 민원 분석 &
실시간 안전 한계 감시 + 정기 보전 \\
\midrule
데이터 민감도 &
금융 거래, 신용 정보 &
건강/유전 정보 (최고 민감) &
행정 데이터, 시민 개인정보 &
산업 비밀, OT 데이터 \\
\bottomrule
\end{tabularx}
\end{table}


\subsection{산업별 거버넌스 자동 구성 도구}
\label{sec:5.5.2}

산업 특성에 따라 거버넌스 정책을 자동으로 구성하는 도구는
조직이 빠르게 적절한 거버넌스 체계를 수립하는 데 도움을 줄 수 있다.
다음 코드는 산업 유형과 AI 활용 사례를 입력받아
맞춤형 거버넌스 구성을 생성하는 Python 클래스이다.

\begin{lstlisting}[language=Python, caption={Industry-Specific AI Governance Configurator}, label={lst:governance_configurator}]
from dataclasses import dataclass, field
from typing import Dict, List, Optional
from enum import Enum

class IndustryType(Enum):
    FINANCE = "finance"
    HEALTHCARE = "healthcare"
    PUBLIC_SECTOR = "public_sector"
    MANUFACTURING = "manufacturing"

class RiskLevel(Enum):
    LOW = "low"
    MEDIUM = "medium"
    HIGH = "high"
    CRITICAL = "critical"

@dataclass
class GovernancePolicy:
    """Governance Policy for a Specific Industry and Use Case"""
    industry: IndustryType
    use_case: str
    risk_level: RiskLevel
    required_approvals: List[str]
    monitoring_frequency: str
    explainability_requirements: List[str]
    human_oversight_type: str
    data_protection_measures: List[str]
    audit_requirements: List[str]
    genai_specific_controls: List[str]

class IndustryGovernanceConfigurator:
    """
    Configures governance policy based on industry type.

    Generates tailored governance requirements considering
    industry regulations, risk profiles, and GenAI considerations.
    """

    # Industry-specific governance templates
    INDUSTRY_PROFILES = {
        IndustryType.FINANCE: {
            "primary_values": ["fairness", "explainability", "consumer_protection"],
            "key_regulations": ["FSC_AI_Guideline", "SR_11-7", "ECOA", "AI_Basic_Act"],
            "default_monitoring": "quarterly",
            "human_oversight": "model_validation_committee",
            "genai_controls": [
                "Hallucination rate monitoring (<2%)",
                "Deepfake detection for KYC/AML",
                "Prompt filtering for PII",
                "RAG-based response verification",
                "AI agent action audit trail",
                "Pre-approval for all GenAI use cases",
            ],
        },
        IndustryType.HEALTHCARE: {
            "primary_values": ["patient_safety", "clinical_validity", "health_equity"],
            "key_regulations": ["MFDS_Guideline", "Medical_Device_Act", "FDA_GMLP"],
            "default_monitoring": "continuous",
            "human_oversight": "clinician_final_decision",
            "genai_controls": [
                "Clinical information accuracy validation",
                "Patient consent for AI-generated advice",
                "Prohibition of unsupervised diagnosis",
                "Medical terminology hallucination check",
                "Subgroup performance equity monitoring",
            ],
        },
        IndustryType.PUBLIC_SECTOR: {
            "primary_values": ["democratic_values", "due_process", "transparency"],
            "key_regulations": ["Intelligence_Info_Act", "Admin_Procedure_Act", "PIPA"],
            "default_monitoring": "annual_report",
            "human_oversight": "official_final_approval",
            "genai_controls": [
                "AI-generated document labeling",
                "Citizen notification of AI use",
                "Sensitive admin data input prohibition",
                "Chatbot accuracy threshold (>98%)",
                "Accessibility for vulnerable populations",
                "Version control for AI-drafted documents",
            ],
        },
        IndustryType.MANUFACTURING: {
            "primary_values": ["physical_safety", "reliability", "real_time"],
            "key_regulations": ["ISO_10218", "SDPA", "IEC_62443"],
            "default_monitoring": "real_time",
            "human_oversight": "supervisor_approval_with_e_stop",
            "genai_controls": [
                "Process parameter generation safety bounds",
                "Operational envelope validation",
                "Fail-safe mode for GenAI suggestions",
                "Latency constraints for real-time control",
            ],
        },
    }

    RISK_ESCALATION = {
        RiskLevel.LOW: {
            "approvals": ["team_lead"],
            "audit_cycle": "annual",
            "explainability": ["global_feature_importance"],
        },
        RiskLevel.MEDIUM: {
            "approvals": ["team_lead", "governance_officer"],
            "audit_cycle": "semi_annual",
            "explainability": ["global_feature_importance", "local_explanations"],
        },
        RiskLevel.HIGH: {
            "approvals": ["team_lead", "governance_officer", "legal_review"],
            "audit_cycle": "quarterly",
            "explainability": [
                "global_feature_importance", "local_explanations",
                "counterfactual_analysis", "adverse_action_notice"
            ],
        },
        RiskLevel.CRITICAL: {
            "approvals": [
                "team_lead", "governance_officer",
                "legal_review", "board_approval"
            ],
            "audit_cycle": "monthly",
            "explainability": [
                "global_feature_importance", "local_explanations",
                "counterfactual_analysis", "adverse_action_notice",
                "full_model_documentation"
            ],
        },
    }

    def configure(
        self,
        industry: IndustryType,
        use_case: str,
        risk_level: RiskLevel,
        uses_genai: bool = False,
    ) -> GovernancePolicy:
        """Generate governance policy for given industry and use case."""

        profile = self.INDUSTRY_PROFILES[industry]
        risk_config = self.RISK_ESCALATION[risk_level]

        # Data protection based on industry
        data_measures = self._get_data_protection(industry)

        # GenAI-specific controls
        genai_controls = profile["genai_controls"] if uses_genai else []

        # Monitoring: use stricter of risk-based and industry default
        monitoring = self._resolve_monitoring(
            risk_config["audit_cycle"], profile["default_monitoring"]
        )

        return GovernancePolicy(
            industry=industry,
            use_case=use_case,
            risk_level=risk_level,
            required_approvals=risk_config["approvals"],
            monitoring_frequency=monitoring,
            explainability_requirements=risk_config["explainability"],
            human_oversight_type=profile["human_oversight"],
            data_protection_measures=data_measures,
            audit_requirements=[
                f"Audit cycle: {risk_config['audit_cycle']}",
                f"Key regulations: {', '.join(profile['key_regulations'])}",
                "Independent validation required"
                    if risk_level in (RiskLevel.HIGH, RiskLevel.CRITICAL)
                    else "Self-assessment acceptable",
            ],
            genai_specific_controls=genai_controls,
        )

    def _get_data_protection(self, industry: IndustryType) -> List[str]:
        measures = {
            IndustryType.FINANCE: [
                "Credit information encryption",
                "Transaction data anonymization",
                "5-year record retention",
            ],
            IndustryType.HEALTHCARE: [
                "PHI de-identification (Safe Harbor)",
                "IRB approval for research use",
                "Pseudonymization per PIPA",
            ],
            IndustryType.PUBLIC_SECTOR: [
                "Administrative data anonymization",
                "Citizen consent management",
                "Data minimization principle",
            ],
            IndustryType.MANUFACTURING: [
                "OT/IT network segmentation",
                "Trade secret protection",
                "Sensor data integrity verification",
            ],
        }
        return measures.get(industry, [])

    def _resolve_monitoring(self, risk_freq: str, industry_freq: str) -> str:
        priority = ["real_time", "monthly", "quarterly", "semi_annual", "annual", "annual_report"]
        risk_idx = priority.index(risk_freq) if risk_freq in priority else 5
        ind_idx = priority.index(industry_freq) if industry_freq in priority else 5
        return priority[min(risk_idx, ind_idx)]

    def generate_summary(self, policy: GovernancePolicy) -> str:
        """Generate human-readable governance policy summary."""
        summary = f"""
=== AI Governance Policy Summary ===
Industry: {policy.industry.value}
Use Case: {policy.use_case}
Risk Level: {policy.risk_level.value}

[Approvals Required]
{chr(10).join(f'  - {a}' for a in policy.required_approvals)}

[Monitoring]: {policy.monitoring_frequency}
[Human Oversight]: {policy.human_oversight_type}

[Explainability Requirements]
{chr(10).join(f'  - {e}' for e in policy.explainability_requirements)}

[Data Protection]
{chr(10).join(f'  - {d}' for d in policy.data_protection_measures)}

[Audit Requirements]
{chr(10).join(f'  - {a}' for a in policy.audit_requirements)}
"""
        if policy.genai_specific_controls:
            summary += f"""
[GenAI-Specific Controls]
{chr(10).join(f'  - {g}' for g in policy.genai_specific_controls)}
"""
        return summary


# Usage Example
configurator = IndustryGovernanceConfigurator()

# Finance: LLM-based customer chatbot (high risk, uses GenAI)
finance_policy = configurator.configure(
    industry=IndustryType.FINANCE,
    use_case="LLM-based customer service chatbot",
    risk_level=RiskLevel.HIGH,
    uses_genai=True,
)
print(configurator.generate_summary(finance_policy))

# Public Sector: AI-assisted document drafting (medium risk, uses GenAI)
public_policy = configurator.configure(
    industry=IndustryType.PUBLIC_SECTOR,
    use_case="AI-assisted administrative document drafting",
    risk_level=RiskLevel.MEDIUM,
    uses_genai=True,
)
print(configurator.generate_summary(public_policy))
\end{lstlisting}


\subsection{교차 산업 거버넌스 평가 도구}
\label{sec:5.5.3}

\begin{practicaltool}{교차 산업 AI 거버넌스 성숙도 평가표}
\textbf{목표}: 조직의 AI 거버넌스 수준을 산업별 핵심 요건 기준으로 평가하고,
개선 우선순위를 도출한다.

\textbf{사용 방법}: 각 항목에 대해 현재 수준을 1점(미도입)부터
5점(최적화)까지 평가하시오.

\begin{table}[htbp]
\centering
\begin{tabularx}{\textwidth}{|p{0.5cm}|p{5cm}|X|p{1.2cm}|}
\hline
\textbf{No.} & \textbf{평가 항목} & \textbf{산업별 가중치 (금/의/공/제)} & \textbf{점수} \\
\hline
\multicolumn{4}{|c|}{\textbf{A. 공정성 및 편향 관리}} \\
\hline
A1 & 보호 속성별 성능 차이 정기 분석 & 높음/중간/높음/낮음 & /5 \\
\hline
A2 & 편향 완화 조치 적용 및 효과 검증 & 높음/중간/높음/낮음 & /5 \\
\hline
A3 & 공정성 지표의 정량적 기준 설정 & 높음/중간/중간/낮음 & /5 \\
\hline
\multicolumn{4}{|c|}{\textbf{B. 설명가능성 및 투명성}} \\
\hline
B1 & 개별 결정에 대한 설명 생성 능력 & 높음/높음/높음/중간 & /5 \\
\hline
B2 & 이해관계자 맞춤형 설명 제공 & 높음/높음/높음/중간 & /5 \\
\hline
B3 & AI 활용 사실 공개 및 고지 체계 & 중간/높음/높음/낮음 & /5 \\
\hline
\multicolumn{4}{|c|}{\textbf{C. 안전성 및 신뢰성}} \\
\hline
C1 & 실시간 성능 모니터링 체계 & 중간/높음/낮음/높음 & /5 \\
\hline
C2 & Fail-safe 메커니즘 구현 & 낮음/높음/낮음/높음 & /5 \\
\hline
C3 & 이상 상황 탐지 및 대응 절차 & 높음/높음/중간/높음 & /5 \\
\hline
\multicolumn{4}{|c|}{\textbf{D. 인간 감독 및 책임성}} \\
\hline
D1 & 인간 감독 체계(HITL) 구축 & 높음/높음/높음/높음 & /5 \\
\hline
D2 & 책임 소재 명확화 및 문서화 & 높음/높음/높음/높음 & /5 \\
\hline
D3 & 이의 신청/불복/구제 절차 마련 & 높음/중간/높음/중간 & /5 \\
\hline
\multicolumn{4}{|c|}{\textbf{E. 생성형 AI 거버넌스}} \\
\hline
E1 & 환각(Hallucination) 모니터링 체계 & 높음/높음/높음/중간 & /5 \\
\hline
E2 & LLM 입출력 기록 및 보존 & 높음/중간/높음/낮음 & /5 \\
\hline
E3 & AI 생성 콘텐츠 표시 및 추적 & 높음/높음/높음/중간 & /5 \\
\hline
E4 & 딥페이크/AI 사기 대응 체계 & 높음/중간/중간/낮음 & /5 \\
\hline
\multicolumn{4}{|c|}{\textbf{F. 데이터 보호 및 규제 준수}} \\
\hline
F1 & 산업별 데이터 보호 규정 준수 & 높음/높음/높음/중간 & /5 \\
\hline
F2 & 모델 문서화 및 기록 보존 & 높음/높음/중간/중간 & /5 \\
\hline
F3 & 외부 감사 대응 체계 & 높음/높음/중간/중간 & /5 \\
\hline
\end{tabularx}
\end{table}

\textbf{점수 해석 기준}:
\begin{itemize}
 \item \textbf{85점 이상 (20문항 $\times$ 5점 = 100점 만점)}:
 선도 수준---산업 Best Practice를 충족하며 벤치마크 역할 가능
 \item \textbf{65--84점}:
 성숙 수준---핵심 요건을 충족하나 일부 영역 개선 필요
 \item \textbf{45--64점}:
 발전 수준---기본 체계는 갖추었으나 체계적 보완 필요
 \item \textbf{44점 이하}:
 초기 수준---거버넌스 체계 수립이 시급하며, 산업별 가중치가 ``높음''인 항목부터 우선 개선
\end{itemize}

\textbf{활용 팁}: 산업별 가중치가 ``높음''인 항목에서 3점 미만을 받은 경우,
해당 영역은 규제 위반 또는 중대한 리스크 노출 가능성이 있으므로
즉각적인 개선 계획을 수립해야 한다.
\end{practicaltool}


\subsection{핵심 시사점}
\label{sec:5.5.4}

본 장의 분석을 통해 도출된 핵심 시사점은 다음과 같다:

\begin{enumerate}
 \item \textbf{공통 기반 위의 산업 특화}:
 모든 산업에 공통적으로 적용되는 거버넌스 기반(데이터 보호, 인간 감독,
 문서화)이 존재하지만, 각 산업의 규제 환경과 리스크 프로파일에 따라
 우선순위와 구현 방식이 달라져야 한다.

 \item \textbf{생성형 AI의 범산업적 도전}:
 LLM과 생성형 AI의 확산으로 환각 관리, 딥페이크 대응,
 AI 생성 콘텐츠 표시 등 새로운 거버넌스 과제가 모든 산업에
 공통적으로 부상하고 있다. 다만 금융의 딥페이크 사기,
 공공의 행정문서 정확성, 의료의 임상 정보 환각 등
 산업별로 가장 치명적인 리스크 유형은 다르다.

 \item \textbf{규제 수렴과 상호 학습}:
 AI 기본법 시행령, EU AI Act 등 범산업적 규제가 시행됨에 따라,
 산업 간 거버넌스 요구사항이 수렴하는 경향이 있다.
 한 산업에서 발전시킨 Best Practice(예: 금융의 4/5 규칙,
 의료의 하위 집단 성능 분석)를 다른 산업에 적용할 수 있다.

 \item \textbf{실무 도구의 중요성}:
 원칙적 선언만으로는 거버넌스가 실현되지 않으며,
 산업별 체크리스트, 자동화된 평가 도구, 투명성 보고서 생성기 등
 구체적인 실무 도구가 뒷받침되어야 한다.
\end{enumerate}

조직은 자신이 속한 산업의 특수성을 이해하고, 본 장에서 제시한 산업별 체크리스트와 실무 도구를 활용하여 규제 준수를 넘어선 비즈니스 가치를 창출해야 한다. 이어지는 제6장에서는 이러한 거버넌스 원칙과 프로세스를 기술적으로 뒷받침하는 \textbf{AI 거버넌스 도구 및 기술}에 대해 상세히 다룬다.

% NEXT_SECTION_HERE
